% Source printed pp. 108-113. Subsections 4.1.5-4.1.6.
\subsection{\texorpdfstring{$\mathcal M_g$}{Mg} 的代数性}\label{src:4.1.5-subsection}
为什么 $\mathcal M_g$ 是代数的? 一个理由是: 每条光滑, 连通的射影曲线 $C$ 都由甚丰沛线丛 $\Omega_C^{\otimes3}$ 给出三重典范嵌入 $C\hookrightarrow\mathbb P^{5g-6}$, 而参数化三重典范嵌入的光滑曲线族的局部闭子概形 $H'\subseteq\Hilb^P(\mathbb P^{5g-6})$ 给出光滑呈示 $H'\to\mathcal M_g$. 不过, 技术细节稍显复杂. $\cBun_{r,d}(C)$ 的代数性 (\bookref{4.1.23}) 会更容易证明.
\begin{theorem}[$\mathcal M_g$ 的代数性]\label{src:4.1.17}
若 $g\geq2$, 则 $\mathcal M_g$ 是 $\Spec\mathbb Z$ 上的代数叠. 此外, $\mathcal M_g\cong[H'/\PGL_{5g-5}]$, 其中 $H'$ 是某个射影 Hilbert 概形的局部闭子概形.
\end{theorem}
\begin{proof}
与证明 $\mathcal M_g$ 是叠 (命题~\bookref{3.5.14}) 时一样, 我们使用光滑曲线族的性质 (\bookref{6.1.16}): 对光滑曲线族 $p:D\to S$, $\Omega_{D/S}^{\otimes3}$ 在 $S$ 上相对甚丰沛, $p_*(\Omega_{D/S}^{\otimes3})$ 是秩为 $5(g-1)$ 的向量丛. 因而 $\Omega_{D/S}^{\otimes3}$ 定义 $S$ 上的闭浸入 $D\hookrightarrow\mathbb P(p_*(\Omega_{D/S}^{\otimes3}))$.
由 Riemann--Roch 定理 (\bookref{6.1.2}), 纤维 $D_s\hookrightarrow\mathbb P^{5g-6}_{\kappa(s)}$ 的 Hilbert 多项式为
\[
P(n):=\chi(\mathcal O_{D_s}(n))=\deg(\Omega_{D_s}^{\otimes3n})+1-g=(6n-1)(g-1).
\]
我们定义
\[
H:=\Hilb^P(\mathbb P^{5g-6}_{\mathbb Z})
\]
为参数化 $\mathbb P^{5g-6}$ 中 Hilbert 多项式为 $P$ 的闭子概形的 (射影) Hilbert 概形 (定理~\bookref{2.1.2}). 设 $C\hookrightarrow\mathbb P^{5g-6}\times H$ 是万有闭子概形, $\pi:C\to H$ 是投影.
我们断言存在局部闭子概形 $H'\subset H$, 使得族 $C_{H'}\to H'$ 满足:
\begin{enumerate}[(a)]
\item 对每个 $h\in H'$, $C_h\to\Spec\kappa(h)$ 光滑且几何连通;
\item 由闭浸入 $C_{H'}\hookrightarrow\mathbb P^{5g-6}\times H'$ 诱导的自然映射
$p_{2,*}\mathcal O_{\mathbb P^{5g-6}\times H'}(1)\to p_{2,*}\mathcal O_{C_{H'}}(1)$
是同构 (等价地, 对所有 $h\in H'$, $H^0(\mathbb P^{5g-6}_{\kappa(h)},\mathcal O(1))\to H^0(C_h,\mathcal O_{C_h}(1))$ 是同构);
\item 线丛 $\Omega_{C_{H'}/H'}^{\otimes3}$ 与 $\mathcal O_{C_{H'}}(1)$ 相差一个从 $H'$ 拉回的线丛.
\end{enumerate}
此外, 若概形态射 $T\to H$ 使得族 $C_T\to T$ 满足 (a)--(c), 则 $T\to H$ 通过 $H'$ 分解. 特别地, $H'\subseteq H$ 唯一. 注意 (a)--(c) 蕴含对 $h\in H$, $C_h\hookrightarrow\mathbb P^{5g-6}_{\kappa(h)}$ 由完备线性系 $\Omega_{C_h/\kappa(h)}^{\otimes3}$ 嵌入.

固有态射 (有限表现) 的纤维光滑是在靶上的开条件 (推论~\bookref{A.3.9}), 因而 $H$ 上 $C_h$ 光滑的条件是开的.
考虑 Stein 分解~\cite[推论 $11.5$]{Har77}
$C\to\widetilde H=\Spec_H\pi_*\mathcal O_C\to H$, 其中 $C\to\widetilde H$ 的纤维几何连通, $\widetilde H\to H$ 有限. 由于 $\mathcal O_H\to\pi_*\mathcal O_C$ 的核与余核具有闭支集, $\widetilde H\to H$ 在 $H$ 的某个开子概形上是同构, 而这恰好就是 $C\to H$ 的纤维几何连通的地方.
综上, 满足 (a) 的点 $h\in H$ 组成 $H$ 的开子概形, 记为 $H_1$.
为使 (b) 成立, 注意上同调与基变换 (\bookref{A.6.8}) 蕴含 $p_{2,*}\mathcal O_{\mathbb P^{5g-6}\times H_1}(1)$ 与 $p_{2,*}\mathcal O_{C_{H_1}}(1)$ 都是向量丛, 其构造与基变换可交换. Riemann--Roch 定理 (\bookref{6.1.2}) 进一步说明两者秩相同.
因此, $\alpha:p_{2,*}\mathcal O_{\mathbb P^{5g-6}\times H_1}(1)\to p_{2,*}\mathcal O_{C_{H_1}}(1)$ 是同构当且仅当它满射. 令 $H_2:=H_1\setminus\operatorname{Supp}(\operatorname{coker}(\alpha))$, 则 (a)--(b) 在 $H_2$ 上成立.
显然, 使得 (a)--(b) 成立的任意映射 $T\to H$ 均通过 $H_2$ 分解.
为使性质 (c) 成立, 由于 $C_{H_2}\to H_2$ 的纤维几何整且相对典范层 $\Omega:=\Omega_{C_{H_2}/H_2}$ 是线丛, 可以使用命题~\bookref{A.6.18}~得到闭子概形 $H_3\hookrightarrow H_2$, 使得态射 $T\to H_2$ 通过 $H_3$ 分解当且仅当 $\Omega^{\otimes3}|_{C_T}$ 与 $\mathcal O_C(1)|_{C_T}$ 相差一个从 $T$ 拉回的线丛.
子概形 $H':=H_3$ 满足 (a)--(c) 以及上述泛性质.

$\mathbb Z$ 上的群概形 $\PGL_{5g-5}=\Aut(\mathbb P^{5g-6}_{\mathbb Z})$ 自然作用于 $H=\Hilb^P(\mathbb P^{5g-6}_{\mathbb Z})$: 若 $g\in\PGL_{5g-5}$ 且 $[D\subseteq\mathbb P^{5g-6}_S]\in H(S)$, 则
$g\cdot[D\subseteq\mathbb P^{5g-6}_S]=[g(D)\subseteq\mathbb P^{5g-6}_S]$.
闭子概形 $H'\subseteq H$ 是 $\PGL_{5g-5}$-不变的.
我们断言 $\mathcal M_g\cong[H'/\PGL_{5g-5}]$. 由商叠的代数性 (\bookref{4.1.10}), 这一断言即可完成定理的证明.
为验证断言, 考虑由 Hilbert 概形的万有族的限制 $C_{H'}\to H'$ 定义的态射 $H'\to\mathcal M_g$. 此态射忘掉嵌入, 即将闭子概形 $D\subseteq\mathbb P^{5g-6}_S$ 送到族 $D\to S$.
该态射是 $\PGL_{5g-5}$-不变的, 并下降为预叠态射 $[H'/\PGL_{5g-5}]^{\mathrm{pre}}\to\mathcal M_g$.
先证明这个映射完全忠实: 设 $S\to H'$ 对应于闭子概形
\[
\begin{tikzcd}
D\arrow[r,hook]\arrow[dr,"q"']&\mathbb P^{5g-6}\times S\arrow[d,"p_2"]\\
&S.
\end{tikzcd}
\]
我们使用下列等价:
\begin{align*}
H^0(\mathbb P^{5g-6}_{\mathbb Z},\mathcal O(1))\otimes\mathcal O_S
&\cong p_{2,*}\mathcal O_{\mathbb P^{5g-6}\times S}(1)\\
&\cong q_*\mathcal O_D(1)&&\text{(性质 (c))}\\
&\cong q_*(\Omega_{D/S}^{\otimes3}\otimes q^*M)&&\text{(性质 (b), $M\in\Pic(S)$)}\\
&\cong q_*(\Omega_{D/S}^{\otimes3})\otimes M&&\text{(投影公式)}.
\end{align*}
$D\to S$ 的自同构诱导 $\Omega_{D/S}^{\otimes3}$ 的唯一自同构, 从而诱导 $q_*(\Omega_{D/S}^{\otimes3})\otimes M$ 的唯一自同构, 后者进而诱导保持 $D$ 的 $\mathbb P^{5g-6}\times S$ 的唯一自同构.
这说明自然映射 $\Stab_{\PGL_{5g-5}}([D\subseteq\mathbb P^{5g-6}_S])\xrightarrow{\sim}\Aut(D\to S)$ 是双射.

由于 $\mathcal M_g$ 是叠 (命题~\bookref{3.5.14}), 叠化的泛性质给出态射 $[H'/\PGL_{5g-5}]\to\mathcal M_g$; 由于 $[H'/\PGL_{5g-5}]^{\mathrm{pre}}\to[H'/\PGL_{5g-5}]$ 完全忠实 (习题~\bookref{3.5.21}), 该映射完全忠实.
只需再验证 $[H'/\PGL_{5g-5}]\to\mathcal M_g$ 本质满.
对此, 只需验证若 $q:D\to S$ 是光滑曲线族, 则存在平展覆盖 $\{S_i\to S\}$, 使得每个 $D_{S_i}$ 都在 $H'(S_i)\to\mathcal M_g(S_i)$ 的像中.
由于 $\Omega_{D/S}^{\otimes3}$ 定义 $S$ 上的闭浸入 $D\hookrightarrow\mathbb P(q_*(\Omega_{D/S}^{\otimes3}))$, 且 $q_*(\Omega_{D/S}^{\otimes3})$ 局部自由, 秩为 $5g-5$, 可以直接取 Zariski 开覆盖 $\{S_i\}$ (因而也是平展覆盖), 使得每个 $(q_*(\Omega_{D/S})^{\otimes3})|_{S_i}$ 都自由.
于是 $\mathcal M_g\cong[H'/\PGL_{5g-5}]$, 而定理~\bookref{4.1.10}~说明 $\mathcal M_g$ 是代数的.
\end{proof}
\begin{remark}\label{src:4.1.18}
所有光滑曲线组成的整个叠 $\mathcal M$ (例~\bookref{3.4.10}) 也是代数的, 因为 $\mathcal M=\coprod_g\mathcal M_g$.
\end{remark}
\begin{exercise}[有益的练习]\label{src:4.1.19}
通过对每个概形 $T$ 直接证明 $\mathcal M_g(T)\cong[H'/\PGL_{5g-5}](T)$, 给出定理~\bookref{4.1.17}~中等价 $\mathcal M_g\cong[H'/\PGL_{5g-5}]$ 的另一证明.

\textit{提示:} Hilbert 概形的万有族给出映射 $[H'/\PGL_{5g-5}](T)\to\mathcal M_g(T)$.
为构造逆映射, 对曲线族 $\pi:C\to T$, 考虑秩为 $5g-5$ 的 (可能非平凡的!) 向量丛 $\pi_*(\Omega_{C/T}^{\otimes3})$ 及诱导的 $T$ 上的闭浸入 $C\hookrightarrow\mathbb P(E)$.
使用射影标架丛 $q:\operatorname{PFr}(E)\to T$ (见习题~\bookref{B.1.57}) 和同构
$\mathcal O_{\operatorname{PFr}(E)}^{\oplus5g-5}\xrightarrow{\sim}q^*E\otimes\mathcal O_{\operatorname{PFr}(E)}(1)$,
得到闭浸入 $C\times_T\operatorname{PFr}(E)\hookrightarrow\mathbb P(q^*E)\cong\mathbb P^{5g-6}\times\operatorname{PFr}(E)$, 以及诱导的 $\PGL_{5g-5}$-等变态射 $\operatorname{PFr}(E)\to H'$.
\end{exercise}
亏格 $g$ 的 $n$-标点光滑曲线族是带有 $n$ 个截面 $\sigma_1,\ldots,\sigma_n:S\to X$ 的概形光滑固有态射 $X\to S$, 使得对每个 $s:\Spec k\to S$, $X_s$ 是亏格 $g$ 的曲线, 且 $\sigma_1(s),\ldots,\sigma_n(s)\in X_s$ 互不相同.
记 $\mathcal M_{g,n}$ 为参数化亏格 $g$ 的 $n$-标点光滑曲线族的 $\Sch_{\text{ét}}$ 上的预叠. 椭圆曲线的情形 $\mathcal M_{1,1}$ 已在习题~\bookref{3.5.15}~中讨论.
\begin{exercise}\label{src:4.1.20}
证明若 $2g-2+n>0$, 则 $\mathcal M_{g,n}$ 是 $\mathbb Z$ 上的代数叠, 即 $g\geq2$, 或 $g=1$ 且 $n\geq1$, 或 $g=0$ 且 $n\geq3$.
\end{exercise}
\begin{exercise}[椭圆曲线的模]\label{src:4.1.21}
\begin{enumerate}[(a)]
\item 利用 Weierstrass 形式 $y^2=x^3+ax+b$ (见~\cite[§3.1]{Sil09}), 证明将素数 $2,3$ 变为可逆元后, 有同构
\[
\mathcal M_{1,1}\times_{\mathbb Z}\mathbb Z[1/6]\cong[(\mathbb A^2\setminus V(\Delta))/\mathbb G_m],
\]
其中作用为 $t\cdot(a,b)=(t^4a,t^6b)$, $\Delta$ 是判别式 $4a^3+27b^2$. 另见~\cite[§3, 定理 $1$]{MS72}.
\item 定义域 $k$ 上的稳定椭圆曲线为二元组 $(E,p)$, 其中 $E$ 是 $k$ 上算术亏格为 $1$ 的不可约射影曲线, 奇点至多为结点, 且 $p\in E(k)$ 是光滑点.
概形 $S$ 上的稳定椭圆曲线族是固有, 平坦且有限表现的态射 $E\to S$, 连同截面 $\sigma:S\to E$, 使得每个纤维都是稳定椭圆曲线.
记 $\overline{\mathcal M}_{1,1}$ 为分类稳定椭圆曲线的 $\Sch$ 上的叠, 证明
\[
\overline{\mathcal M}_{1,1}\times_{\mathbb Z}\mathbb Z[1/6]\cong[(\mathbb A^2\setminus0)/\mathbb G_m],
\]
其中作用与上面相同.
\end{enumerate}
\end{exercise}
\begin{exercise}\label{src:4.1.22}
回顾 $\mathcal M_{0,0}\cong B\PGL_2$ (习题~\bookref{3.5.15}). 证明:
\begin{enumerate}[(a)]
\item $\mathcal M_{0,1}\cong BU_2$, 其中 $U_2\subseteq\PGL_2$ 是由上三角矩阵组成的二维子群;
\item $\mathcal M_{0,2}\cong B\mathbb G_m$;
\item $\mathcal M_{0,3}\cong\Spec\mathbb Z$;
\item 当 $n>3$ 时, $\mathcal M_{0,n}\cong(\mathbb P^1\setminus\{0,1,\infty\})^{n-3}\setminus\Delta$, 其中 $\Delta$ 是 $n-3$ 个点中至少有两个相等的闭子概形.
\end{enumerate}
\end{exercise}

\subsection{\texorpdfstring{$\cBun(C)$}{Bun(C)} 的代数性}\label{src:4.1.6-subsection}
命题~\bookref{3.5.17}~证明了, 对域 $k$ 上的概形 $X$, $\QCoh(X)$, $\Coh(X)$ 与 $\cBun(X)$ 都是 $(\Sch/k)_{\text{ét}}$ 上的叠. 现在专门考虑曲线 $C$ 的情形, 尽管下面的结果在远为一般的情形下也成立 (见习题~\bookref{4.1.25}).
我们将
\[
\Coh_{r,d}(C)\subseteq\Coh(C)\qquad\text{与}\qquad\cBun_{r,d}(C)\subset\cBun(C)
\]
定义为参数化二元组 $(E,S)$ 的全子范畴, 要求对每个几何点 $\Spec K\to S$, $E_K$ 都是 $C_K$ 上秩为 $r$, 次数为 $d$ 的凝聚层.
\begin{theorem}[$\cBun(C)$ 的代数性]\label{src:4.1.23}
若 $C$ 是代数闭域 $k$ 上的光滑, 连通的射影曲线, 则叠 $\cBun(C)$ 与 $\Coh(C)$ 是代数的, 且 $\cBun(C)\subseteq\Coh(C)$ 是开子叠.
对整数 $r\geq0$ 与 $d$, $\cBun_{r,d}(C)$ 与 $\Coh_{r,d}(C)$ 是代数叠, 而 $\cBun_{r,d}(C)\subseteq\cBun(C)$ 与 $\Coh_{r,d}(C)\subseteq\Coh(C)$ 是开闭子叠.
\end{theorem}
\begin{remark}\label{src:4.1.24}
定理给出分解
\[
\Coh(C)=\coprod_{r,d}\Coh_{r,d}(C)\qquad\text{与}\qquad\cBun(C)=\coprod_{r,d}\cBun_{r,d}(C).
\]
尽管 $\cBun_{r,d}(C)$ 与 $\Coh_{r,d}(C)$ 不拟紧 (定义~\bookref{4.3.21}), 下面的证明表明, $\cBun_{r,d}(C)$ 或 $\Coh_{r,d}(C)$ 的每个拟紧开子叠, 都是拟射影概形 (事实上是 Quot 概形的开子概形) 关于 $\GL_N$ 的商叠.
\end{remark}
\begin{proof}
首先注意, $\Coh_{r,d}(C)$ 与 $\cBun_{r,d}(C)$ 是 $(\Sch/k)_{\text{ét}}$ 上的叠, 因为它们由几何纤维上的条件定义为叠 $\Coh(C)$ 与 $\cBun(C)$ 的全子范畴.
为说明 $\cBun(C)\subseteq\Coh(C)$ 是开子叠, 设态射 $S\to\Coh(C)$ 由 $C\times S$ 上的有限表现的拟凝聚层 $E$ 分类, 且 $E$ 在 $S$ 上平坦.
若 $V\subseteq C\times S$ 是 $E$ 为向量丛的开区域, 则 $S\setminus p_2(C\times S\setminus V)$ 是开的, 并与纤维积 $S\times_{\Coh(C)}\cBun(C)$ 等同.
由凝聚层的 Riemann--Roch 定理 (\bookref{9.1.2}), $C$ 上秩为 $r$, 次数为 $d$ 的凝聚层 $E$ 的 Hilbert 多项式为
\[
P(n):=\chi(E(n))=\deg(E(n))+\operatorname{rk}(E(n))(1-g)=d+rn+r(1-g).
\]
由于 Hilbert 多项式在平坦族中局部常值 (命题~\bookref{A.2.4}), 包含 $\cBun_{r,d}(C)\hookrightarrow\cBun(C)$ 与 $\Coh_{r,d}(C)\hookrightarrow\Coh(C)$ 是开闭子叠. 因而只需证明 $\Coh_{r,d}(C)$ 是代数的.

若 $E$ 是 $C$ 上秩为 $r$, 次数为 $d$ 的凝聚层, 则 Serre 消没定理说明, 对 $N\gg0$, $E(N)$ 由整体截面生成且 $H^1(C,E(N))=0$.
特别地, $\Gamma(C,E(N))\otimes_k\mathcal O_C\twoheadrightarrow E(N)$ 满射, 且在整体截面上诱导同构.
对每个整数 $N$, 记
\[
\mathcal U_N\subseteq\Coh_{r,d}(C)
\]
为参数化二元组 $(S,E)\in\Coh_{r,d}(C)$ 的子叠, 要求对每个 $s\in S$, $E_s(N):=E(N)|_{C\times\Spec\kappa(s)}$ 在 $C_s:=C\times_k\kappa(s)$ 上由整体截面生成, 且 $H^1(C_s,E_s(N))=0$.
这些条件蕴含 $P(N)=h^0(C,E_s(N))$, $R^1p_{2,*}E(N)=0$, 以及 $p_2^*p_{2,*}E(N)\to E(N)$ 满射.
我们断言 $\mathcal U_N\subseteq\Coh_{r,d}(C)$ 是开子叠.
为验证断言, 设 $S$ 是概形, $E$ 是 $C\times S$ 上在 $S$ 上平坦且 Hilbert 多项式为 $P$ 的有限表现的拟凝聚层, 考虑图
\[
\begin{tikzcd}
&C\times S\arrow[dl,"p_1"']\arrow[dr,"p_2"]&\\
C&&S.
\end{tikzcd}
\]
由上同调与基变换的一次直接应用 (见命题~\bookref{A.6.12}), 使得 $H^1(C,E_s(N))=0$ 的点 $s\in S$ 组成的区域 $S'\subseteq S$ 是开的.
此外, $(R^1p_{2,*}E(N))|_{S'}=0$, 而 $(p_{2,*}E(N))|_{S'}$ 是秩为 $P(N)$ 的向量丛, 其构造与基变换可交换.
层 $F:=\operatorname{coker}(p_2^*p_{2,*}E(N)\to E(N))$ 具有闭支集, 开子集 $S'':=S'\setminus p_2(\operatorname{Supp}(F))\subseteq S$ 是满足 $E_s(N)$ 由整体截面生成且 $H^1(C,E_s(N))=0$ 的点 $s\in S$ 组成的区域.
若 $S\to\Coh_{r,d}(C)$ 是分类 $E$ 的映射, 则基变换 $\mathcal U_N\times_{\Coh_{r,d}(C)}S$ 与 $S''$ 等同, 断言得证.

对每个 $N$, 考虑 Quot 概形
\[
Q_N:=\Quot^P(\mathcal O_C(-N)^{\oplus P(N)}/C/k),
\]
它参数化 Hilbert 多项式为 $P$ 的商 $\mathcal O_C(-N)^{\oplus P(N)}\twoheadrightarrow F$ (定理~\bookref{2.1.3}).
与上面类似的论证说明, 存在开子概形 $Q'_N\subseteq Q_N$, 参数化满足
$H^0(q(N)):H^0(C,\mathcal O_C)^{\oplus P(N)}\to H^0(C,F(N))$ 满射且 $H^1(C,F(N))=0$ 的商 $q:\mathcal O_C(-N)^{\oplus P(N)}\twoheadrightarrow F$.
Quot 概形 $Q_N$ 从 $\GL_{P(N)}$ 继承自然作用: 给定 $g\in\GL_{P(N)}$ 与 $[q:\mathcal O_C(-N)^{\oplus P(N)}\twoheadrightarrow F]\in Q_N$, 令
\[
g\cdot q:=
\bigl[\mathcal O_C(-N)^{\oplus P(N)}\xrightarrow{g^{-1}}
\mathcal O_C(-N)^{\oplus P(N)}\xrightarrow{q}F\bigr]\in Q_N.
\]
区域 $Q'_N\subseteq Q_N$ 显然是 $\GL_{P(N)}$-不变的. 态射
\[
Q'_N\longrightarrow\mathcal U_N,\qquad
[\mathcal O_C(-N)^{\oplus P(N)}\twoheadrightarrow F]\longmapsto F,
\]
是 $\GL_{P(N)}$-不变的 (即对于 $\mathcal U_N$ 上的平凡作用为 $\GL_{P(N)}$-等变的), 因而下降为预叠态射
\[
\Psi^{\mathrm{pre}}:[Q'_N/\GL_{P(N)}]^{\mathrm{pre}}\longrightarrow\mathcal U_N.
\]
为说明 $\Psi^{\mathrm{pre}}$ 完全忠实, 考虑商 $[q:\mathcal O_C(-N)^{\oplus P(N)}\twoheadrightarrow F]\in Q_N(k)$ 及诱导映射 $\Stab_{\GL_{P(N)}}(q)\to\Aut(F)$.
自同构 $\alpha\in\Aut(F)$ 诱导 $H^0(C,F(N))\cong k^{P(N)}$ 的自同构, 保持 $H^0(q(N))$ 的核.
这定义了 $\Stab(q)\to\Aut(F)$ 的逆映射, 而将这一构造推广到商的族并不困难.
由于 $\Coh_{r,d}(C)$ 是叠 (命题~\bookref{3.5.17}), $\mathcal U_N$ 也是叠, 并有诱导的叠态射 $\Psi:[Q'_N/\GL_{P(N)}]\to\mathcal U_N$, 它完全忠实 (由习题~\bookref{3.5.21}) 且本质满 (由构造).
于是 $\mathcal U_N=[Q'_N/\GL_{P(m}]$, 且
\[
\Coh_{r,d}(C)=\bigcup_N[Q'_N/\GL_{P(N)}].
\]
商叠的代数性 (定理~\bookref{4.1.10}) 蕴含 $\Coh_{r,d}(C)$ 的代数性.
\end{proof}
\begin{exercise}\label{src:4.1.25}
\begin{enumerate}[(1)]
\item (简单) 修改上述论证, 证明若 $X$ 是域 $k$ 上任意维数的射影概形, 则 $\cBun(X)$ 与 $\Coh(X)$ 是代数叠.
\item (中等) 更一般地, 证明若 $X\to S$ 是 Noether 概形的强射影态射 (见 \S\bookref{2.4.4-subsection}), 则叠 $\Coh(X/S)$ 是代数叠, 它在 $S$-概形 $T$ 上的对象是 $X_T$ 上在 $T$ 上平坦的有限表现的拟凝聚层. (另见~\cite[定理 $4.6.2.1$]{LMB00}.)
\end{enumerate}
另见定理~\bookref{C.7.7}~中使用 Artin 公理的证明.
\end{exercise}
\begin{exercise}[困难]\label{src:4.1.26}
设 $X$ 是域 $k$ 上的射影概形, $G$ 是 $k$ 上的仿射代数群. 定义 $X$ 上主 $G$-丛的叠 $\cBun_G(X)$, 并证明它是代数的.

\textit{提示:} 选择忠实嵌入 $G\hookrightarrow\GL_n$, 并证明主丛的诱导定义了可表态射 $\cBun_G(X)\to\cBun_{\GL_n}(X)$, 其靶为 $X$ 上的向量丛的叠.
\end{exercise}
